\section*{Capítulo \ref{ch:b3:facies-models} --- Modelos de facies}

\begin{solution}{exo:b3:facies-models:1}
Desde el mar hacia tierra: mar abierto, transición de mar abierto con arenas de tormenta,
\omterm{def:b3:facies-models:shoreface}{anteplaya}, playa, una llanura costera con pantanos. Cada ambiente quedó sobre su vecino
situado hacia el mar: la línea de costa se desplazó mar adentro, una regresión. Una transgresión los
habría apilado en orden inverso, con la llanura costera en la base y el lodo de mar abierto en el techo,
a menudo con una superficie de erosión donde pasó la
\omterm{def:b3:facies-models:shoreface}{anteplaya}.
\end{solution}

\begin{solution}{exo:b3:facies-models:2}
(a) Relleno de canal o \omterm{def:b3:facies-models:rivers}{barra de meandro}: una base erosiva con cantos rodados, estratificación cruzada,
después rizos y lodo con raíces. (b) \omterm{def:b3:facies-models:shoreface}{Anteplaya} o frente deltaico en progradación:
lodo con capas de arena de tormenta y \omterm{def:b3:facies-models:hcs}{estratificación cruzada hummocky} debajo,
arena de \omterm{def:b3:facies-models:shoreface}{anteplaya} con estratificación cruzada y playa encima. (c) Arenas de \omterm{def:b3:facies-models:rivers}{río trenzado} apiladas,
una duna eólica o un banco carbonatado: conjuntos de estratificación cruzada de tamaño de grano constante,
con grandes láminas frontales de duna en el caso eólico, o de granos
esqueléticos y oolíticos en el banco.
\end{solution}

\begin{solution}{exo:b3:facies-models:3}
$\phi = (2.65 - 2.40)/(2.65 - 1.0) = 0.152$; con salmuera,
$0.25/1.55 = 0.161$.
\end{solution}

\begin{solution}{exo:b3:facies-models:4}
$e^{-2\pi z/\lambda} = e^{-\pi} = 0.043$ en $z = \lambda/2$: unos
\qty{38}{m} y \qty{55}{m}. A \qty{20}{m}: $e^{-2\pi\times20/75} = 0.19$
y $e^{-2\pi\times20/110} = 0.32$ del movimiento en superficie.
\end{solution}

\begin{solution}{exo:b3:facies-models:5}
$(60 - 47.6)/(189 - 47.6) = 0.088$; $(60 - 52.6)/(189 - 52.6) = 0.054$. Un
valor de la matriz incierto en \qty{5}{\micro\second/ft} modifica la porosidad en
un tercio: la ecuación de tiempo medio necesita un valor de la matriz comprobado en testigos.
\end{solution}

\begin{solution}{exo:b3:facies-models:6}
$(6.0 + 3.5 + 8.0 + 2.5)/40 = 20/40 = 0.50$. Volumen de poros: $20\times0.22
\times10^6 = \qty{4.4e6}{m^3}$ por kilómetro cuadrado.
\end{solution}

\begin{solution}{exo:b3:facies-models:7}
Cerca del cañón: solo Ta, con las divisiones superiores eliminadas por el flujo siguiente; hacia abajo
por el abanico: Tc (rizos) y Td (limo) bajo Te (lodo). La corriente pierde velocidad
y su carga gruesa con la distancia; solo la fracción fina recorre
\qty{30}{km}, y se deposita por un flujo más lento que forma rizos, no
por precipitación desde una suspensión densa.
\end{solution}

\begin{solution}{exo:b3:facies-models:8}
Una \omterm{def:b3:facies-models:rivers}{barra de meandro} de un \omterm{def:b3:facies-models:rivers}{río meandriforme}: depósito residual sobre el fondo erosionado del canal,
arena con estratificación cruzada del canal profundo, arena con rizos en la parte alta de la barra,
lodo de llanura de inundación con raíces. La sucesión registra el canal desde el fondo hasta
la orilla, de modo que el canal tuvo unos \qty{9}{m} de profundidad (algo menos, tras
la compactación del lodo). No a través de la base: es una superficie de erosión, y
lo que yacía debajo pudo ser cualquier depósito más antiguo.
\end{solution}

\begin{solution}{exo:b3:facies-models:9}
$60/40\,000 = 0.0015$, \ang{0.09}: una rampa. $300/1500 = 0.20$, \ang{11.3},
tras una cima plana de \qty{5}{m} de profundidad: una \omterm{def:b3:facies-models:ramp}{plataforma bordeada}, con granos de arrecife o de banco
en el margen, y bloques y flujos granulares en la pendiente inferior.
\end{solution}

\begin{solution}{exo:b3:facies-models:10}
La franja tarda un tiempo $w/v$ en pasar por un punto fijo, durante el cual el fondo
se eleva respecto al nivel del mar a razón de $A$ bajo ella: la arena llena $h = A\,w/v$.
$2000\times0.2/20 = \qty{20}{m}$. Con $A = 0$ la \omterm{def:b3:facies-models:shoreface}{anteplaya} deja solo
el espesor de su propio perfil, una capa delgada, pues las franjas se deslizan unas sobre otras
sin agradar. La razón $A/v$ es la tangente del ángulo de trayectoria
de la proposición del capítulo~\ref{ch:b3:sequence-stratigraphy}: el
espesor de cada facies es su anchura multiplicada por esa tangente.
\end{solution}

\begin{solution}{exo:b3:facies-models:11}
Densidad: $(2.65 - 2.30)/1.65 = 0.212$; sónico: $(86 - 52.6)/136.4 = 0.245$.
Con $\phi = 0.245$, $\rho_{\mathrm{ma}} = (2.30 - 0.245)/(1 - 0.245) =
\qty{2.72}{g/cm^3}$: más denso que el cuarzo, como corresponde a una arenisca cementada por
calcita (2.71) o con granos de calcita.
\end{solution}

\begin{solution}{exo:b3:facies-models:12}
El lodo de mar abierto sobre un carbón pone en contacto ambientes no contiguos: una
\omterm{def:b3:sequence-stratigraphy:parasequence}{superficie de inundación}, aquí una \omterm{def:b3:sequence-stratigraphy:surfaces}{superficie transgresiva} de erosión por oleaje, con los restos de conchas
como depósito residual. Al desplazarse la línea de costa tierra adentro, el oleaje rebajó el fondo
hasta aproximadamente el \omterm{def:b3:facies-models:shoreface}{nivel de base del oleaje de buen tiempo}: la playa, la parte alta de la \omterm{def:b3:facies-models:shoreface}{anteplaya} y el
techo de la llanura costera, hasta unos \qty{10}{m} de columna, pueden
faltar.
\end{solution}

\begin{solution}{pb:b3:facies-models:1}
\textbf{1.} Creta (fría, de baja densidad), lutita desde \qty{766}{m} (caliente), caliza desde
\qty{866}{m} (fría, densa, rápida), arenisca desde \qty{1046}{m} (formas acampanadas),
esquisto desde \qty{1166}{m}.

\textbf{2.} Creta \qty{766}{m}, lutita 100, caliza 180, arenisca
\qty{120}{m}; base de la cobertera en $441.0 - 1166 = \qty{-725}{m}$.

\textbf{3.} $(2.71 - 2.03)/1.71 = 0.40$ cerca de la superficie, $(2.71 -
2.36)/1.71 = 0.20$ a \qty{750}{m}: la creta pierde la mitad de su porosidad
bajo el peso de la carga suprayacente, por compactación.

\textbf{4.} Sus granos son arcillas de densidad incierta, y su agua está
en parte ligada a ellas: fallan las hipótesis de la ecuación.

\textbf{5.} Basamento cristalino: se atraviesa la cobertera y no se espera
encontrar ningún yacimiento por debajo.

\textbf{6.} $(2.71 - 2.573)/1.71 = 0.080$.

\textbf{7.} $(58.9 - 47.6)/141.4 = 0.080$: lo mismo.

\textbf{8.} $(2.65 - 2.32)/1.65 = 0.200$.

\textbf{9.} $(79.9 - 52.6)/136.4 = 0.200$: lo mismo.

\textbf{10.} $0.33/1.55 = 0.213$, un punto y un tercio más.

\textbf{11.} La arenisca, con \qty{20}{\percent} de porosidad en sus arenas;
la lutita sobre la caliza la sella, y los lodos de llanura de inundación dividen la
arenisca en compartimentos.

\textbf{12.} Una campana: rayos gamma bajos en la arena con estratificación cruzada, que ascienden
a lo largo del intervalo granodecreciente hasta los valores altos del lodo con raíces, sobre una base neta.

\textbf{13.} Una \omterm{def:b3:facies-models:association}{sucesión granodecreciente}: lag del fondo del canal sobre una base erosiva,
dunas de canal (estratificación cruzada), techo de la \omterm{def:b3:facies-models:rivers}{barra de meandro} (rizos), llanura de
inundación (lodo con raíces).

\textbf{14.} El canal, su \omterm{def:b3:facies-models:rivers}{barra de meandro} en la orilla interna, diques naturales y una
llanura de inundación con suelos, uno junto a otro.

\textbf{15.} Aproximadamente el espesor de la arena y del intervalo granodecreciente, $10.4 + 3.2 =
\qty{13.6}{m}$, algo más antes de la compactación.

\textbf{16.} Un \omterm{def:b3:facies-models:rivers}{río meandriforme} cuyo cinturón de canales se desplazó repetidamente por
\omterm{def:b3:facies-models:rivers}{avulsión} mientras la cuenca subsidía, cortando cada nuevo canal
la llanura de inundación que dejó el anterior.

\textbf{17.} No: cada base es erosiva; el lodo de llanura de inundación situado debajo y la
arena de canal situada encima no tienen por qué haber sido vecinos.

\textbf{18.} $(60 - 30)/(105 - 30) = 0.40$.

\textbf{19.} Mínimo: $14\times(0.60 + 0.40\times0.25) = \qty{9.80}{m}$;
máximo: $16\times(0.65 + 0.40\times0.20) = \qty{11.68}{m}$.

\textbf{20.} Arena limpia \qty{58.8}{m}; intervalos granodecrecientes \qty{33.4}{m}; lodo
\qty{27.8}{m}.

\textbf{21.} Neto $= 58.8 + 0.40\times33.4 = \qty{72.2}{m}$; \omterm{def:b3:facies-models:ntg}{razón neto--bruto}
$72.2/120 = 0.60$.

\textbf{22.} Fracciones 0.20 y 0.60: \qty{65.5}{m} y \qty{78.8}{m}, una
dispersión de \qty{13}{m} debida solo a la elección del límite de corte.

\textbf{23.} $72.2\times0.20\times10^6 = \qty{1.44e7}{m^3}$ por
kilómetro cuadrado.

\textbf{24.} El espesor neto de la arenisca triásica en Varden-2 es
\textbf{\qty[text-series-to-math=true]{72.2}{m}}.
\end{solution}
